\section{Checks of the connection, invariants, and classification}\label{appB}
\subsection{Objects to be checked}
The symbolic calculation constructs the Levi--Civita connection from the
structure coefficients, adds the torsion ansatz, and computes curvature
and contractions. Verification compares exterior-form and component
calculations, rather than merely redisplaying the same expressions.
The scope is three geometries, three torsion modes (AX, VT, MX),
and three action variants (FULL, TT, REE).

Basic conditions are $C^a{}_{bc}=-C^a{}_{cb}$, $d^2e^a=0$,
torsion reconstruction by the Cartan equation, and metric compatibility.
With the last connection index denoting differentiation direction,
\begin{equation}
\nabla_c\delta_{ab}
=-\Gamma^d{}_{ac}\delta_{db}
-\Gamma^d{}_{bc}\delta_{ad}=0.
\tag{B-01}\label{eq:B-01}
\end{equation}
This is equivalent to $\Gamma_{abc}+\Gamma_{bac}=0$.
The geometric torsion represented by an input array is determined
separately by the Cartan equation.

\subsection{Curvature and identities}
Metric-compatible curvature is antisymmetric in the first index pair,
and the two-form definition gives antisymmetry in the last pair.
Pair exchange $R_{abcd}=R_{cdab}$ is not required for a general
connection with torsion. For example, the MX ansatz can give
\begin{equation}
R_{0123}-R_{2301}=\frac{2V\eta}{3s}.
\tag{B-02}\label{eq:B-02}
\end{equation}
This supplies a counterexample to treating pair exchange as a universal
test. The first and second Bianchi identities are checked in their
torsionful exterior-form expressions,
\begin{equation}
DT^a=R^a{}_b\wedge e^b,\qquad DR^a{}_b=0.
\tag{B-03}\label{eq:B-03}
\end{equation}
For NY, construct $d(e^a\wedge T_a)$, $T^a\wedge T_a$, and
$e^a\wedge e^b\wedge R_{ab}$ independently and compare the coefficients
and difference in \eqref{eq:M2-06}.
The cuts $V=0$, $\eta=0$, and $\theta=0$ also check the agreement of
action variants and the EC part.

\subsection{Symbolic verdicts and numerical comparisons}
The engine distinguishes \texttt{PROVED\_ZERO}, where zero is proved;
\texttt{WITNESS\_NONZERO}, where a nonzero witness is found; and
\texttt{UNPROVED}, where neither is established.
An unproved expression is not treated as zero, and a concrete nonzero
example is distinguished from proof of a symbolic identity.

The reduced potentials can be checked by independently calculating
curvature and volume, integrating, and comparing with
Section~\ref{sec3}. Further independent calculations cover first and
second derivatives, discriminants, positive stationary roots, and
action differences between maxima and minima.
Representative cases include minima near the origin, large-radius
minima, $C=0$, $D=0$, $V=0$, and both torsion coefficients zero.

Agreement at finitely many numerical points does not replace an identity.
The analytic formulas are checked symbolically, while independent
multiprecision evaluation checks classifier roots, heights, and
near-boundary behavior.
Comparisons require corresponding $T$ and $\tau$ parameters, the adopted
volume, and the same variation path.

\subsection{Reproducing calculations}
Record the code version, inputs, logs, and outputs together.
For each formula comparison, specify the geometry, action variant,
torsion mode, and identity tested.
Appendix~\ref{appF} gives execution procedures; the
\href{https://github.com/Muacca/DPPUv2-paper01/blob/main/script/docs/CONVENTIONS_ja.md}{component conventions}
provide the code's geometric definitions.

\subsection{Figures and summary tables}
See Table~\ref{tab:v3:5} for the conditions and values.

\begin{table}[htbp]
\centering
\small
\caption{Local invariants at $s=1,V=-4,\eta=3$. Here $\mathcal T^2=T_{abc}T^{abc}$; $n$ is the density relative to the volume form, and $n_{\rm FULL}=n_{\rm TT}-n_{\rm REE}$.}\label{tab:v3:5}
\setlength{\tabcolsep}{4pt}
\renewcommand{\arraystretch}{1.25}
\begin{tabular}{lrrrrr}
\toprule
Space & $\mathcal R_{\rm LC}$ & $\mathcal R_{\rm EC}$ & $\mathcal T^2$ & $n_{\rm TT}=n_{\rm REE}$ & $n_{\rm FULL}$ \\
\midrule
$S^3$ & $24$ & $-122/3$ & $680/3$ & $48$ & $0$ \\
$T^3$ & $0$ & $-194/3$ & $680/3$ & $48$ & $0$ \\
$\mathrm{Nil}^3$ & $-1/2$ & $-391/6$ & $680/3$ & $48$ & $0$ \\
\bottomrule
\end{tabular}
\end{table}

\FloatBarrier
